\section{五连鞭}
\begin{enumerate}
  \item 总体说明。
        所谓法治五连鞭是指对法律运行过程中立法、执法、司法、守法和法律监督这 5 大环节的描述。
        这 5 个部分可用性极强，
        一方面每一个部分都可以单独命题考查，
        另一方面不管是法治思想整体的落实，
        还是某一个坚持的落实，
        都可以从这 5 个方面来展开，
        因为任何原则或理念的落地，
        永远是要体现在这 5 个具体的法治工作当中的，
        所以 5 连鞭可以理解是\bs{万能“怎么办”}。
        为保持逻辑的一致性 5 连鞭中每一个部分，
        依然按照总体价值要求和具体措施 2 个部分展开，
        大家要尽可能精细的掌握。
        以下几点作为特别提示，大家要格外注意:
        \begin{enumerate}
          \item 法治思想中对具体法治工作的强调只有前 4 个部分，
                但出于知识完整性的考虑，
                我增加了监督环节，
                答题语言全部来自于官方文件，
                可靠性毋庸置疑，
                可以放心背通。
                另外，从过去的考试情况来看，
                确实也出现过“先上车，后补票”的情况，
                因此对法律监督的掌握也能起到预防效果。
          \item 5 连鞭中每一个部分的价值意义，
                不要忘记\bs{万能价值}，
                这也可以在相当程度上减负，
                但最好是把这里背的和之前背的万能价值结合起来。
          \item 一般来说，5 连鞭在整体应用时，把第一段压缩一下，足够解决问题，
                但 5 连鞭的每一个具体部分都可能具体考查（比如去年），
                所以我在压缩的基础上尽可能给大家全面的展现知识。
                当然，考虑到大家的背诵压力，
                在此提示 2 点首先\bs{具体措施}部分，
                只要\bs{记住关键句}即可，
                展开部分可以按照字面意思发挥或者进行\bs{正反论证}，
                搞一搞废话文学；
                其次，具体措施是为了应对“怎么办”这一部分的答题要求，
                如果实在没有具体展开的能力，
                可以\bs{回到宏观层面}发挥，
                也就是借用之前背通过的\bs{政治方向（前 3 个坚持，党、人、路）和重要保障（法律人和少数人）} 2 个部分，
                道理很简单，
                不管是立法，
                执法还是任何的法治工作，
                \bs{强调坚持正确的政治方向，
                  重视人才培养，
                  充分发挥关键少数的领导职能，
                  都是没有任何问题的}。
                总之，可以虚实结合！
          \item 5 连鞭有没有重点呢，
                当然有，
                法治社会部分，
                也就是守法这一环节，是绝对的，超级重点，
                无论是“我国不能成为诉讼大国”，
                还是“枫桥经验”等新提法和新指示，
                都直接关联这个部分，
                所以要求法治社会，
                全部背诵！
                有困难，克服即可，这个没得商量！
        \end{enumerate}
  \item 科学立法。
        \begin{enumerate}
          \item 总体意义和要求。
                \begin{enumerate}
                  \item \bs{法律是治国之重器，良法是善治之前提}。
                  \item 建法系，必须坚持\zx{立法先行}，
                  \item 发挥立法的\zx{引领和推动}作用，抓住提高\zx{立法质量}这个关键。
                  \item 要恪守\zx{以民为本、立法为民}理念，\zx{贯彻社会主义核心价值观}，
                  \item 使每一项立法都\zx{符合宪法精神、反映人民意志、得到人民拥护}。
                  \item 要把\bs{公正、公平、公开}原则\zx{贯穿到}立法的全过程，
                  \item \zx{完善}立法体制机制，坚持\zx{立、改、废、释}并举，
                  \item \zd{增强}法律法规的\zx{及时性、系统性、针对性、有效性}。
                \end{enumerate}
          \item 具体措施。
                \begin{enumerate}
                  \item 健全\zx{宪法实施和监督}制度。
                  \item 完善\zx{立法体制}。深入推进\zx{科学立法、民主立法}，加强\zx{重点领域立法}。
                \end{enumerate}
        \end{enumerate}
  \item 严格执法（法治政府）。
        \begin{enumerate}
          \item 总体意义和要求。
                \begin{enumerate}
                  \item 执法是行政机关\zx{履行政府职能}、\zx{管理经济社会事务}的主要方式。
                  \item \bs{法律的生命力在于实施，法律的权威也在于实施}。
                  \item 各级政府必须坚持\zx{在党的领导下}、\zx{在法轨上}开展工作，
                  \item 创新\zx{执法体制}，完善执法\zx{程序}，严格执法\zx{责任}，推进\zx{综合执法}，
                  \item 建立\bs{权责统一、高效权威}的依法行政体制，
                  \item 加快建设\bs{权责法定、执法严明、公开公正、廉洁高效、守法诚信}的法治政府。
                \end{enumerate}
          \item 具体措施。
                \begin{enumerate}
                  \item 依法\zx{全面履行}政府职能。
                  \item 健全\zx{依法决策}机制。
                  \item 深化行政\zx{执法体制}改革。
                  \item 坚持严格规范公正\zx{文明执法}。
                  \item 强化对行政权力的\zx{制约和监督}。
                  \item 全面推进\zx{政务公开}。
                \end{enumerate}
        \end{enumerate}
  \item 公正司法（去年考过）。
        \begin{enumerate}
          \item 总体意义和要求。
                \begin{enumerate}
                  \item 公正是法治的\bs{生命线}，司法公正对社会公正具有重要的\zx{引领}作用，司法不公对社会公正具有致命的\zx{破坏}作用。
                  \item 英国的哲学家\zx{培根}曾经有云，一次不公正的审判的危害结果要远胜于十次犯罪，因为后者不过是污染了水流，而前者却是污总染了水源。
                  \item 必须完善\zx{司法管理体制和司法权力运行机制}，推进\zx{司法体制改革}，
                  \item \zx{规范司法行为}，加强对\zx{司法活动的监督}，确保司法机关\zx{独立行使}职权，
                  \item 保障\zx{案件处理}的过程符合\bs{程序正义}，案件处理的结果符合\bs{实体正义}，
                  \item 努力让人民群众在每一个司法案件中\bs{感受到}公平正义。
                \end{enumerate}
          \item 具体措施。
                \begin{enumerate}
                  \item 完善、优化司法职权配置，确保\zx{依法、独立、公正行使审判权和检察权}。
                  \item 推进\zx{严格司法}，\bs{坚持以事实为根据，以法律为准绳}，
                  \item 健全\bs{[事实认定符合客观真相、办案过程符合程序公正、办案结果符合实体公正]}的法律制度。
                  \item 保障\bs{人民群众参与司法}，
                        坚持\zx{人民司法为人民}，依靠\zx{人民推进公正司法}，通过公正司法\zx{维护人民权益}。
                  \item 加强\zx{人权司法}保障。
                  \item 加强对\zx{司法活动的监督}。
                \end{enumerate}
        \end{enumerate}
  \item 全民守法（也就是法治社会，这是今年的超级重点）！
        \begin{enumerate}
          \item 总体意义和要求。
                \begin{enumerate}
                  \item 法治社会是构筑法治国家的\zx{基础}，全民依法治国需要全社会\zx{共同参与}。
                  \item \bs{人民权益要靠法律保障，法律权威要靠人民维护}。
                  \item 卢梭曾经说过：“一切法律之中最重要的法律，既不是铭刻在\zx{大理石}上，也不是刻在\zx{铜表}上，而是铭刻在\zx{公民的内心}里”。
                  \item 法律的权威源自\zx{人民的内心拥护和真诚信仰}，法律\zx{要发挥}功能，需要全社会\zx{信仰法律}。
                  \item 必须\zx{弘扬}社会主义\zx{法治精神}，\zx{建设}社会主义\zx{法治文化}，增强全社会\zx{厉行法治的积极性和主动性}。
                  \item 形成\bs{守法光荣、违法可耻}的社会氛围，
                  \item 使全体人民都成为\zx{社法治忠实的崇尚者、自觉的遵守者、坚定的捍卫者}。
                \end{enumerate}
          \item 具体措施。
                人民是依法治国的\bs{主体和力量源泉}，
                三治的建设构成\zx{有机联系的系统整体}，
                法治事业的顺利发展，\zx{离不开全民}的参与和支持。
                \bs{在今后的法治社会建设过程中，必须要坚持以下基本要求}。
                \begin{enumerate}
                  \item 推动全社会\zx{树立法治意识}。加强\zx{普法教育和普法宣传}，引导民众\zx{树立正确的}法治意识和法治观念。
                  \item 推进\zx{多层次多领域依法治理}。综合利用各种手段，广泛发挥\zx{社会各方主体在矛盾化解，社会关系协调}方面的积极价值。
                  \item 建设\zx{完备的法律服务体系}。推进覆盖城乡居民的\zx{公共法律服务体系}建设，加强\zx{民生领域法律服务}。
                  \item 健全\zx{依法维权和纠纷化解}机制。引导民众\zx{合法合理表达利益诉求}，将\zx{社会矛盾的解决纳入法治轨道}。
                \end{enumerate}
        \end{enumerate}
  \item 法治监督（预防性补充）。
        \begin{enumerate}
          \item 总体意义和要求。
                \begin{enumerate}
                  \item 党和国家监督体系是党在\zx{长期执政}条件下\bs{实现自我净化、自我完善、自我革新、自我提高}的重要制度保障。
                  \item \bs{严密的法律监督体系}，是进一步\zx{完善}\bs{社会主义民主政治制度}的必然要求，
                  \item 是全面\zx{贯彻落实}依法治国战略，\zx{保障}立执司守带修饰的应有之义，
                  \item 也是确保\zx{公权力}始终\zx{在法轨上运行}，\zx{保障人民合法权益}的必然内涵。
                  \item 必须健全党统一领导、全面覆盖、权威高效的\zx{监督体系}，
                  \item 增强\zx{监督严肃性、协同性、有效性}，
                  \item 形成决策科学、执行坚决、监督有力的\zx{权力运行机制}，
                  \item 确保\zx{赋权在轨谋福、利}。
                \end{enumerate}
          \item 具体措施。
                \begin{enumerate}
                  \item \zx{健全党和国家监督制度}。
                  \item \zx{加强对法律实施的监督}。
                \end{enumerate}
        \end{enumerate}
\end{enumerate}